% Do not change document class, margins, fonts, etc.
\documentclass[a4paper,oneside,bibliography=totoc]{scrartcl}

% some useful packages (add more as needed)
\usepackage[utf8]{inputenc}
\usepackage{graphicx}
\usepackage{latexsym}
\usepackage{amsmath}
\usepackage{amssymb}
\usepackage{tabularx}
\usepackage{booktabs}
\usepackage{algorithm} % you can modify the algorithm style to your liking
\usepackage{algorithmic}
\usepackage{csquotes}
\renewcommand{\algorithmiccomment}[1]{\hfill\textit{// #1}}
\usepackage[usenames,dvipsnames]{xcolor}
\usepackage[colorlinks,citecolor=Green]{hyperref} % you may change/remove the colors
\usepackage{lipsum} % you do not need this
\usepackage{environ}
\usepackage{tikz}
\usetikzlibrary{positioning,shapes.gates.logic.US,shapes.gates.logic.IEC,automata,arrows, backgrounds}
\usepackage{natbib}
\bibliographystyle{chicagoa}
\setcitestyle{authoryear,round,semicolon,aysep={},yysep={,}} \let\cite\citep

% example enviroments and commands (add more as needed)
\newtheorem{definition}{Definition} \newtheorem{proposition}{Proposition}

\usepackage{enumitem}
\setlist[itemize]{topsep=.15em,itemsep=.15em,parsep=0pt,leftmargin=1.3em}

% images are looked up in the main folder first, then in figures/
\graphicspath{{figures/}}

% figure + caption inside an answer box (height capped so the box cannot overflow)
\newcommand\answerfig[2]{%
  {\centering\includegraphics[width=0.92\linewidth,height=5.5cm,keepaspectratio]{#1}\par}%
  \vspace{.15em}{\footnotesize #2\par}}

\newcommand\todo[1]{\textcolor{red}{\framebox{\footnotesize TODO}\expandafter\ifx\expandafter\relax\detokenize{#1}\relax\else\ #1\fi}}


% DO NOT MODIFY
\newlength\oldparskip
\NewEnviron{answer}[1][8cm]{%
  \sbox{0}{\parbox[t][#1][t]{0.99\linewidth}{%
      \setlength{\oldparskip}{\parskip}%
      \setlength{\parskip}{.5em}%
      \BODY%
      ~%
      \setlength{\parskip}{\oldparskip}%
    }}%

  \begin{tikzpicture}
    % inner sep restores the \fbox padding
    \node[inner sep=\fboxsep, outer sep=0pt] (A) {\usebox{0}};
    % border thickness like \fboxrule
    \draw[line width=\fboxrule] (A.south west) rectangle (A.north east);
  \end{tikzpicture}%
  \vskip\baselineskip
}


\begin{document}

\title{ML26, Assignment 2}
\author{Abhaykumar Narendrakumar Monasara (amonasar, S2188510)\\
  Selma Nezihoglu (snezihog, 2390917)}
\date{\today}

\maketitle

\section*{Task 1a) and Task 1c)}
\begin{answer}[19.5cm]
  \small
  \answerfig{task1_densities.png}{KDEs of all 57 features: (a) raw data (provided plot),
  (b) raw word frequencies up to their 99th percentile, (c) z-scored training data on $[-1,5]$.}

  \textbf{1a)} The provided plot (a) is almost uninformative: all curves collapse into a single spike at $0$.
  The $x$-axis has to cover the largest value, \texttt{capital\_run\_length\_total} (up to $15841$), whereas
  word and character frequencies never exceed $42.8$ and $32.5$, respectively. On one axis spanning such
  different scales, no individual distribution is visible. Plotting the raw word frequencies only up to
  their 99th percentile (b) shows what plot (a) hides:
  \begin{itemize}
    \item \emph{Sparsity.} For the median feature, $84\%$ of the training emails have the value exactly $0$.
    \item \emph{Right skew.} The distributions are strongly right-skewed (median skewness $9.4$): almost all
          mass lies near $0$, followed by a long tail formed by a few emails with large values.
    \item \emph{Exceptions.} Only the three capital-run-length features are never $0$, and they also have by
          far the largest range.
  \end{itemize}

  \textbf{1c)} We z-scored both training and test data using the mean and standard deviation of the
  training data. All features now share one scale, so every density is visible (c), and on the training data
  each feature has mean $0$ and variance $1$. Because z-scoring is an affine map with positive slope, it
  changes only location and scale, not shape: each curve is still a sharp peak with a long right tail. The
  peak lies slightly \emph{below} $0$, since the many zero entries are mapped to $-\mu_j/\sigma_j<0$. Three
  observations stand out:
  \begin{itemize}
    \item \emph{Extreme outliers.} Z-scores range from $-0.93$ to $45.1$, although $98.4\%$ of all values lie
          in $[-1,3]$. The largest values belong to \texttt{capital\_run\_length\_longest} ($45.1$),
          \texttt{capital\_run\_length\_average} ($40.3$) and \texttt{char\_freq\_\#} ($38.1$); in each case a
          single email lies more than $38$ standard deviations above the mean.
    \item \emph{Asymmetry.} Since the raw features are non-negative, the z-scores of feature $j$ are bounded
          below by $-\mu_j/\sigma_j$, and no value falls below $-0.93$. The right tails, by contrast, are very
          long, so the features are far from Gaussian.
    \item \emph{Test data is not exactly standardized.} Test means lie in $[-0.06,0.09]$ and test variances in
          $[0.05,2.76]$. This is expected: the test data is transformed with the training statistics, so no
          information from the test set enters preprocessing.
  \end{itemize}
\end{answer}


\section*{Task 2a)}
\begin{answer}[19.5cm]
  \small
  \textbf{Claim.} Consider $p(y=1\mid\mathbf x)=\sigma(\mathbf w^\top\mathbf x+b)$ with empirical risk
  \[
    \hat R(\mathbf w,b)=\frac1N\sum_{i=1}^N \ell(y_i,s_i),\quad s_i=\mathbf w^\top\mathbf x_i+b,\quad
    \ell(y,s)=-y\log\sigma(s)-(1-y)\log\sigma(-s).
  \]
  Rescale and shift every feature, $\tilde x_{ij}=a_j x_{ij}+c_j$ with constants $a_j\neq0$, $c_j\in\mathbb R$.
  Then the smallest achievable empirical risk (the infimum, which is a minimum whenever a minimizer exists)
  is the same on $\tilde{\mathbf X}$ and on $\mathbf X$.

  \textbf{Proof.} The loss of example $i$ depends on the parameters only through its score $s_i$. Parameters
  that yield the same score for every example therefore have the same empirical risk. Given any
  $(\mathbf w,b)$ for the original data, set
  \[
    \tilde w_j=\frac{w_j}{a_j},\qquad \tilde b=b-\sum_{j=1}^D\frac{w_j c_j}{a_j}.
  \]
  Then for every example $i$
  \[
    \tilde{\mathbf w}^\top\tilde{\mathbf x}_i+\tilde b
    =\sum_{j=1}^D\frac{w_j}{a_j}(a_j x_{ij}+c_j)+b-\sum_{j=1}^D\frac{w_j c_j}{a_j}
    =\mathbf w^\top\mathbf x_i+b,
  \]
  hence $\hat R_{\tilde{\mathbf X}}(\tilde{\mathbf w},\tilde b)=\hat R_{\mathbf X}(\mathbf w,b)$. The map
  $(\mathbf w,b)\mapsto(\tilde{\mathbf w},\tilde b)$ is a bijection with inverse $w_j=a_j\tilde w_j$,
  $b=\tilde b+\sum_j\tilde w_j c_j$. Both problems therefore attain exactly the same set of risk values, and
  their infima coincide. In particular, if $(\mathbf w^*,b^*)$ is an ERM estimate on $\mathbf X$, its image
  $(\tilde{\mathbf w}^*,\tilde b^*)$ is an ERM estimate on $\tilde{\mathbf X}$ with the same risk and the
  same predictions. \hfill$\square$

  The bias term is essential for the shift: without $b$, the term $\sum_j w_j c_j/a_j$ cannot be absorbed.
  Pure rescaling ($c_j=0$) still leaves the minimal risk unchanged, but a shift in general does not.

  \textbf{Why z-scores, then?} The ERM solutions are equivalent, but \emph{finding} them is not equally easy:
  \begin{itemize}
    \item \emph{Optimization.} Gradient descent uses a single step size for all weights. On the raw data the
          features differ by orders of magnitude (Task 1a), so a step size that is stable along the
          large-scale features is far too small for the others. Accordingly, the condition number of
          $\mathbf X^\top\mathbf X$ drops from $1.8\cdot10^{9}$ on the raw data to $2.9\cdot10^{3}$ on the
          z-scores. With the settings of Task 2 (500 epochs, $\varepsilon_0=0.01$, bold driver), GD reaches
          risk $0.2145$ and training accuracy $92.2\%$ on z-scores, but only risk $2.38$ and accuracy
          $49.7\%$ on the raw data.
    \item \emph{No bias.} Our model has no bias term (Task 2b), so the argument above does not cover shifts:
          z-scoring changes which decision functions the model can represent. Centering the features at $0$
          partly compensates for the missing bias.
    \item \emph{Regularization and interpretation.} The $L_2$ penalty of Task 4 is not invariant to
          rescaling, since features on a small scale need large weights and are therefore penalized more.
          Weights can also only be compared (Task 3c) when all features are on the same scale.
  \end{itemize}
\end{answer}

\section*{Task 2e)}
\begin{answer}[19.5cm]
  \small
  \answerfig{task2_gd_sgd.png}{Settings: $\mathbf w_0\sim\mathcal N(\mathbf 0,\mathbf I)$
  (seed 0), $\varepsilon_0=0.01$, 500 epochs, bold driver. $\hat R^*\approx0.2047$ is estimated from a
  GD run with 30\,000 epochs.}

  \resizebox{\linewidth}{!}{%
  \begin{tabular}{lcccccc}
    \toprule
     & $\hat R$ after 500 ep. & epochs to $1.1\,\hat R^*$ & time/epoch & train acc. & test acc. & $\lVert\mathbf w\rVert_2$ \\
    \midrule
    GD  & 0.2145 & 144 & 0.09\,ms & 92.2\% & 91.9\% & 8.9 \\
    SGD & 0.2164 & 8   & 4.6\,ms  & 92.1\% & 91.7\% & 6.8 \\
    \bottomrule
  \end{tabular}}

  \textbf{Per epoch (a).} SGD makes far more progress early on. Its risk is $0.32$ after one epoch, and
  it comes within $10\%$ of $\hat R^*$ after 8 epochs. This is expected, since an SGD epoch performs
  $N=3065$ updates while a GD epoch performs one. GD barely improves over roughly the first 50 epochs
  and needs 144 epochs to reach the same level.

  \textbf{Step sizes (c).} GD's slow start is a step-size problem: with $\varepsilon_0=0.01$, a single
  full-gradient step per epoch moves $\mathbf w$ very little. The bold driver increases the step size
  by $5\%$ per epoch until it reaches about $60$ in epoch 179. After that it oscillates between roughly
  $8$ and $35$ (23 halvings), and GD overtakes SGD. For SGD, the bold driver is counterproductive.
  Gradient noise from single examples often raises the risk from one epoch to the next, which triggers
  36 halvings. The step size decays to about $10^{-3}$, and SGD stalls at a risk of $0.216$. Three
  further seeds give the same final risk to within $\pm0.00002$, so the random example order has a
  negligible effect here.

  \textbf{Per running time (b).} In wall-clock time, GD is clearly faster. A GD epoch consists of
  vectorized matrix--vector products, whereas an SGD epoch is a Python loop over 3065 examples and takes
  about 50 times as long. GD reaches a lower risk after $0.04$\,s than SGD does after $2.3$\,s, i.e.\
  after all 500 epochs. Part of this gap comes from the implementation rather than the algorithms; a
  vectorized or compiled SGD loop would be considerably faster per epoch.

  \textbf{Convergence.} Neither method has converged after 500 epochs, and both remain about $0.01$
  above $\hat R^*$. In the 30\,000-epoch reference run, the risk continues to decrease slowly while
  $\lVert\mathbf w\rVert_2$ grows to about $50$. For prediction, the remaining gap matters little. Both
  models reach about $92\%$ training and test accuracy, and the final GD and SGD weight vectors have a
  correlation of $0.96$.

  \textbf{Conclusion.} On this small data set, full-batch GD is the better choice. Its epochs are cheap,
  and the bold driver lets it exploit large step sizes. SGD becomes worthwhile when the full gradient is
  expensive to compute, i.e.\ for very large $N$. To converge, however, it needs a decreasing step-size
  schedule rather than a heuristic that reacts to risk increases caused by gradient noise.
\end{answer}


\section*{Task 3b) and 3c)}
\begin{answer}[19.5cm]
  % YOUR ANSWER HERE
\end{answer}

\section*{Task 4b)}
\begin{answer}[19.5cm]
  % YOUR ANSWER HERE
\end{answer}


\section*{Task 4c)}
\begin{answer}[19.5cm]
  % YOUR ANSWER HERE
\end{answer}

\newpage
\section*{Ehrenwörtliche Erklärung}

Ich versichere, dass ich die beiliegende Bachelor-, Master-, Seminar-, oder
Projektarbeit ohne Hilfe Dritter und ohne Benutzung anderer als der angegebenen
Quellen und in der untenstehenden Tabelle angegebenen Hilfsmittel angefertigt
und die den benutzten Quellen wörtlich oder inhaltlich entnommenen Stellen als
solche kenntlich gemacht habe. Alle KI-unterstützten Inhalte wurden von mir
geprüft. Für den Inhalt der Arbeit übernehme ich die volle Verantwortung und bin
in der Lage, ihn auf Nachfrage zu erläutern. Diese Arbeit hat in gleicher oder
ähnlicher Form noch keiner Prüfungsbehörde vorgelegen. Ich bin mir bewusst, dass
eine falsche Erklärung rechtliche Folgen haben wird.

% All AI tools that contributed to this work in terms of content or language
% must be listed, including their name (e.g., ChatGPT, Copilot, Claude,
% Consensus AI, ...) , purpose of use (e.g., code generation, text creation,
% translation, research) and where it has been applied (e.g., which part of the
% thesis, code, ...).
%
% Word-processing and reference-management software, spell-checking tools, and
% non-AI-assisted functions of the development environment do not need to be
% listed. Likewise, AI models that are used as the subject of investigation or
% as a method within the work itself (e.g., LLM-as-a-judge in the evaluation)
% and are already mentioned there do not need to be listed.
%
% The AI declaration must include at least one an entry for every graded subtask
% (i.e., where points are specified). The entry should state whether or not you
% used AI to solve parts of the task. Using AI for things such as language
% improvements or API search does not count as "Solved with AI", but things such
% as code gen, writing drafts, or solving math tasks does.
%
% We may grade tasks that you (partially or fully) solved with AI more
% stringently. So make sure that whenever you use AI, the results are correct.
% Our recommendation is to not use AI to solve the tasks, as this undermines the
% learning goals of this course and, ultimately, your own education.
\begin{center}
  \textbf{Declaration of Used AI Tools} \\[.3em]
  \begin{tabularx}{\textwidth}{lclXc}
    \toprule
    Task & Solved with AI? & Tool & Purpose & Useful? \\
    \midrule
    1a & Yes & Claude & Code for further plots and statistics & yes \\
    1b & Yes & Claude & Code generation (\texttt{normalize\_data}) & yes \\
    1c & Yes & Claude & Code for plots and statistics & yes \\
    2a & Yes & Claude & Draft of proof and written answer, supporting experiment & yes \\
    2b & Yes & Claude & Code generation & yes \\
    2c & Yes & Claude & Code generation (GD and progress plots) & yes \\
    2d & Yes & Claude & Code generation & yes \\
    2e & Yes & Claude & Code for experiments and figures & yes \\
    3a \\
    3b \\
    3c \\
    4a \\
    4b \\
    4c \\
    \bottomrule
  \end{tabularx}
\end{center}

\vspace{2cm}
\noindent Unterschrift\\
\noindent Mannheim, den XX.XX.XXXX \hfill


\end{document}
